\documentclass[11pt]{article}
\usepackage[margin=1.05in]{geometry}
\usepackage[T1]{fontenc}
\usepackage{amsmath,amssymb,amsthm}
\usepackage{booktabs}
\usepackage[colorlinks=true,linkcolor=black,citecolor=black,urlcolor={blue!55!black},
            pdftitle={Accumulated Transport Resistance: A Geometry-Only Surrogate for Diffusive Maintenance Burden},
            pdfauthor={Echo-S Research}]{hyperref}
\allowdisplaybreaks

\newtheorem{proposition}{Proposition}[section]
\newtheorem{lemma}[proposition]{Lemma}
\newtheorem{definition}[proposition]{Definition}
\theoremstyle{remark}
\newtheorem{remark}[proposition]{Remark}

% --- notation ---
\newcommand{\R}{\mathbb{R}}
\newcommand{\Om}{\Omega}
\newcommand{\Sig}{\Sigma}
\newcommand{\Gam}{\Gamma}
\newcommand{\Acr}{A_{\mathrm{cross}}}
\newcommand{\Rgeom}{\mathcal{R}_{\mathrm{geom}}}
\newcommand{\dCtube}{\Delta C_{\mathrm{tube}}}
\newcommand{\Qdn}{Q_{\downarrow}}
\newcommand{\RATR}{R_{\mathrm{ATR}}}
\newcommand{\rhou}{\rho_u}
\newcommand{\rhohat}{\hat\rho_{\mathrm{ATR}}}
\newcommand{\Pimaint}{\Pi_{\mathrm{maint}}}
\newcommand{\RG}{R_{G}}
\newcommand{\Qstar}{Q^{\!*}}
\newcommand{\norminf}[1]{\lVert#1\rVert_{\infty}}
\newcommand{\eps}{\varepsilon}

\title{Accumulated Transport Resistance:\\
\large A Geometry-Only Surrogate for Diffusive Maintenance Burden,
and a Phase-A Protocol to Test It}
\author{Echo-S Research}
\date{July 2026}

\begin{document}
\maketitle

\begin{abstract}
Whether an interior point of a living tissue can be maintained is a transport
question, not a distance question: what matters is how much metabolic demand must
be served across each available transport cross-section, integrated along the way
to a supply surface. The companion paper on the mixed-boundary torsion problem
establishes, with proof, that Euclidean distance to the supply set yields
\emph{no} universal lower bound on the transport burden, and that on tapered
``horn'' geometries a distance-only scale can overestimate the true burden by an
unbounded factor. Building on that forced mathematics, we define the
\emph{accumulated transport resistance} (ATR): a demand-weighted, cross-section-
weighted, path-integrated functional whose exact ground truth is the torsion
field~$u$ solving $-\nabla\!\cdot(P\nabla u)=q$ with a Dirichlet supply portion
and reflecting walls. We separate two objects the earlier work sometimes ran
together---a purely geometric resistance under uniform demand and a
demand-weighted concentration drop---and show they coincide only when demand is
uniform. We then organize the whole problem around a single ratio,
$\Qdn(s)/(P(s)\Acr(s))$, the demand routed through each transport cross-section,
which cleanly separates two regimes: volume-graded tapers, where the burden falls
\emph{below} the distance prediction, and demand-loaded bottlenecks, where it
\emph{exceeds} a matched-distance reference (a true-field amplification
$\rhou=1.83$ is exhibited numerically). Because the exact field is not deployable,
the scientific content is a protocol: a Phase-A ``surrogate contest'' on synthetic
geometries with known~$u$, in which a geometry-only ATR surrogate must pass a
gatekeeper---does it track the exact torsion burden?---before any biological
claim is licensed. We give the surrogate an ablation decomposition that isolates
\emph{which} ingredient carries the signal (on a matched-distance family the gain
is the restriction$\,\times\,$demand interaction, not distance, restriction, or
demand alone), a validation package stronger than correlation, and a candidate
eikonal construction whose error we show to be regime-dependent and largest
exactly where it would be used. We close by demoting the free-energy-principle
reading to an optional interpretive layer and offering, in its place, a
measurable one: form responds not to raw distance but to the accumulated
difficulty of maintaining relation across its own geometry. Every displayed
identity and numerical claim is machine-checked by the companion harness.
\end{abstract}

\tableofcontents

\section{Introduction}\label{sec:intro}

A diffusing metabolite entering a tissue from a supply surface---a capillary
wall, a Haversian canal, the outer boundary of a spheroid---must reach every
interior cell that consumes it. The classical readout of whether it can is a
\emph{distance}: Krogh's cylinder model~\cite{krogh1919} fixed a supply radius,
and much of the downstream tradition asks only how far a cell sits from its
nearest source. Distance is attractive because it is cheap: it needs only the
geometry of the supply set, never a transport solve. It is also, as a predictor
of maintenance burden, provably wrong in a precise and recoverable way.

The companion paper of this line~\cite{torsion} studies the object that actually
governs the steady maintenance problem---the mixed-boundary torsion function
\begin{equation}\label{eq:pde}
-\nabla\!\cdot\!\bigl(\mathcal{A}\nabla u\bigr)=q\ \text{ in }\Om,\qquad
u=0\ \text{ on }\Sig,\qquad \mathcal{A}\nabla u\cdot\nu=0\ \text{ on }\Gam,
\end{equation}
with $\Sig\subset\partial\Om$ the supply (Dirichlet) portion, $\Gam$ the
reflecting (no-flux) remainder, $q>0$ a volumetric demand, and $\mathcal{A}$ a
measurable, uniformly elliptic diffusivity. There $u$ is the accumulated cost of
maintaining an interior point against demand under constrained transport, and its
supremum $\norminf{u}$ is the worst-case burden. That paper proves two things we
take as given here. First, distance to $\Sig$ yields no universal lower bound on
$\norminf{u}$: on a family of tapered horns, all three metric forms of the
conjectured bound $qR^2/(2n\,P)$ fail, and the ratio
$P\norminf{u}/(qR^2)$ has infimum $0$. Second---and this is the point most easily
misread---the failure is that distance \emph{overestimates}: as a horn narrows,
the true burden \emph{shrinks}, while a distance-only surrogate keeps charging for
the length. The corrected lesson is not ``narrowing raises resistance'' but
something sharper, which this paper is built around:
\begin{quote}
Transport burden depends on how much maintenance demand must pass through each
available transport cross-section---not on distance or narrowness in isolation.
\end{quote}

The remainder of the paper turns that sentence into a research program.
Section~\ref{sec:forced} recalls exactly what the torsion theorem forces and,
importantly, what it does \emph{not}, and it separates two transport objects that
must never be multiplied by demand twice. Section~\ref{sec:tiers} descends from
the exact field, which is not deployable, through a PDE-derived flux diagnostic to
a geometry-only surrogate, and states the discipline that governs the descent: a
gatekeeper that every biological claim must pass. Section~\ref{sec:regimes}
introduces the organizing ratio and the two regimes it separates.
Section~\ref{sec:ratios} fixes two amplification numbers that the earlier draft
conflated. Section~\ref{sec:phaseA} is the core: the Phase-A surrogate contest,
its ablation decomposition, its model-comparison discipline, and a validation
package that a high correlation coefficient cannot satisfy. The shorter
Sections~\ref{sec:outcome}--\ref{sec:branch} settle the preregistration of
outcomes, the supply-boundary controls, a dimensionless maintenance number, and
the treatment of branch points as competing primary models rather than a
sensitivity footnote. Section~\ref{sec:lcn} places the lacuno-canalicular network
and its controls. Section~\ref{sec:interp} demotes the free-energy reading to an
optional interpretive layer and states the measurable principle that survives.
Sections~\ref{sec:status}--\ref{sec:verify} record honestly what is forced, what
is a modeling choice, and what remains open, and describe the machine
verification. Throughout, ``$P$'' denotes the diffusivity scale $P_{\max}$ of
\eqref{eq:pde}, and we write $\Acr(s)$ for a transport cross-section and
$V(s)=\int_0^s\Acr$ for the volume upstream of arclength~$s$.

\section{What the torsion theorem forces, and what it does not}\label{sec:forced}

\subsection{The comparison field and the direction of the horn result}

The horn analysis of~\cite{torsion} is driven by an explicit \emph{comparison
field}. On a power-law horn of profile $g(x)=\eps h\,(x/h)^{k/2}$ with constant
scalar diffusivity $P$, set
\begin{equation}\label{eq:vfield}
v(x,y)=a\,(h^2-x^2)-b\,\lvert y\rvert^2,\qquad
a=\frac{q}{P\bigl(2+k(n-1)\bigr)},\qquad b=\frac{ak}{2}.
\end{equation}
This field solves $-P\Delta v=q$ identically and---by the choice $b=ak/2$
matched to the profile exponent---has \emph{exactly vanishing flux through the
lateral wall}. It is not the torsion function: $v$ does not vanish on the supply
mouth $\Sig$ and has positive outflux on the end face. What it delivers is a
rigorous one-sided comparison,
\begin{equation}\label{eq:upper}
0\ \le\ u\ \le\ v+b(\eps h)^2\quad\text{a.e.},
\end{equation}
hence a rigorous \emph{upper} bound on $\norminf{u}$. Equation~\eqref{eq:upper}
is a theorem, established in the companion paper by a Stampacchia truncation
against the reflecting wall; we recall it, and do not reprove it, here.

The apex value of the comparison field satisfies a graded identity that is the
mathematical heart of the whole construction:
\begin{equation}\label{eq:graded}
\frac{q}{P}\int_0^h\frac{V(s)}{\Acr(s)}\,ds
=\frac{q\,h^2}{P\bigl(2+k(n-1)\bigr)}=v(0,0),
\qquad\text{since }\ \frac{V(s)}{\Acr(s)}=\frac{s}{1+\tfrac{k(n-1)}{2}}.
\end{equation}
Read \eqref{eq:graded} carefully: the tube quantity on the left \emph{decreases}
as the taper exponent $k$ grows. Numerically it falls from $0.167$ at $k=2$ to
$0.046$ at $k=10$; its derivative in $k$ is negative throughout. The object that
\emph{grows} without bound under narrowing is not the burden but the
distance-only failure factor $qR^2/\norminf{u}$, which climbs from about $6$ to
about $22$ across the same range. The horns are therefore geometries on which a
distance functional and the true burden separate without bound, in the direction
of distance \emph{over}charging. Whether the true apex value $u(0)$ equals the
graded integral \eqref{eq:graded}, rather than merely being bounded by it,
remains open; the exact, path-free, general-domain object is instead the Green
representation
\begin{equation}\label{eq:green}
u(x)=\int_\Om G(x,y)\,q(y)\,dy,
\end{equation}
valid wherever the mixed Green function exists, for almost every $x$, and
pointwise under the additional regularity used in the computational phases. In
the manuscript's general setting $\mathcal{A}$ may be merely measurable and
uniformly elliptic, so \eqref{eq:green} is scoped to that existence and not
claimed beyond it; for the regular synthetic domains of Phase~A the qualification
imposes no operational limit.

\subsection{Two tube objects that must not be multiplied by demand twice}

A unit ambiguity haunted earlier drafts, and closing it is a prerequisite for
everything downstream. There are two genuinely different one-dimensional tube
integrals, and only one of them already contains demand. Define, separately, the
\emph{geometric resistance} under uniform volumetric demand,
\begin{equation}\label{eq:rgeom}
\Rgeom=\int\frac{V(s)}{P(s)\,\Acr(s)}\,ds,\qquad V(s)=\int_0^s\Acr,
\end{equation}
and the \emph{predicted concentration drop} under general demand,
\begin{equation}\label{eq:dctube}
\dCtube=\int\frac{\Qdn(s)}{P(s)\,\Acr(s)}\,ds,\qquad
\Qdn(s)=\int_{\text{downstream of }s} q(y)\,dy.
\end{equation}
The downstream demand $\Qdn(s)$ already carries the metabolic load, so $\dCtube$
\emph{is} the predicted drop, with the dimensions of a concentration. The
geometric resistance $\Rgeom$ carries only geometry. They are related, but only
under uniform demand, where $\Qdn(s)=q\,V(s)$ and hence
\begin{equation}\label{eq:qrel}
\dCtube=q\,\Rgeom\qquad(\text{uniform }q\text{ only}).
\end{equation}
Equation~\eqref{eq:qrel} is the whole content of the unit split: multiplying the
demand-weighted $\dCtube$ by a further factor of $q$ would double-count the load.
Keeping $\Rgeom$ (volume-weighted) and $\dCtube$ (demand-weighted) as distinct
named objects removes the ambiguity at the source.

A canonical normalized burden requires a preregistered normalization $\Qstar$:
$\RG(x)=u(x)/\Qstar$ is meaningful only once $\Qstar$ is fixed. For uniform demand
the clean choice is $\Qstar=q$, giving $\RG=u/q$ with the dimensions of $\Rgeom$;
total-demand and characteristic-rate normalizations are legitimate alternatives,
but a specimen-specific $\Qstar$ chosen after the fact would obscure any
cross-specimen comparison and is excluded unless declared in advance.

\section{From exact field to deployable surrogate: three tiers and a
gatekeeper}\label{sec:tiers}

The exact torsion field $u$ is the ground truth, but computing it requires solving
the full PDE on the true geometry---precisely what a cheap predictor is meant to
avoid. The program therefore lives on three tiers. At the top is $u$ itself. One
step down is a \emph{PDE-derived flux diagnostic}, built from the true current
$J=-P\nabla u$: it records what an ideal reduction of the field would retain, but
it is not deployable because it still needs $u$. At the bottom is a
\emph{geometry-only surrogate} $\RATR$, computable from the segmented geometry
alone. The scientific question is whether the bottom tier can be made to track the
top.

\paragraph{The gatekeeper.} We elevate one step from ``validation'' to a
structural gate. Solving \eqref{eq:pde} on the actual geometry and testing whether
a constructed $\RATR$ tracks the exact $u$ is not a downstream nicety; it is the
gate every biological claim must pass. Until that tracking is demonstrated across
synthetic geometries where $u$ is known, the honest register is
\begin{quote}
\emph{we are testing a computational surrogate inspired by the exact horn
functional,}
\end{quote}
and never ``we have transported the horn theorem into tissue.'' The theorem is
exact on its stated geometry; on an arbitrary network the surrogate is a surrogate
until the gate is passed. Diffusion through irregular tissue does not descend a
single path---it distributes across a field of paths---so any reduction of the
field to flux tubes or watersheds is an additional modeling move, not something
the theorem forces.

\paragraph{Flux geometry and its caveats.} Because $\nabla\!\cdot J=q$, the
current is not divergence-free: throughput varies along a tube, equal-throughput
tubes need not remain equal, stream tubes encounter critical points and
separatrices, a global tiling may fail or require measure-theoretic conventions,
and the streamlines of $J$ run from the interior \emph{toward} $\Sig$, not inward
from it. The flux-tube construction is therefore a diagnostic that establishes
what an ideal reduction would keep; it is not itself cheap. A deployable predictor
must approximate that diagnostic \emph{without} solving the PDE, or else be
presented openly as a reduced summary of a PDE solution.

\paragraph{Eikonal descent is one candidate, and its error is
regime-dependent.} A natural cheap construction takes transport paths to be the
characteristics of the eikonal equation $\lvert\nabla T\rvert=1/P$, i.e.\ minimum-
cost paths to $\Sig$. These are not flux lines. In a controlled test the eikonal
descent points straight along the supply direction ($0^\circ$ from the $-x$ axis,
demand-blind), while the true current $-\nabla u$ curves $10.6^\circ$ toward the
demand mass. That single global figure understates the stakes, because the angular
error between eikonal descent and the true flux \emph{is} the modeling error
incurred whenever an eikonal ATR is used, and it differs sharply between the two
regimes of Section~\ref{sec:regimes}. Measured inside the transport corridor---the
cells where the true flux is dominantly longitudinal toward $\Sig$, about
$20{,}000$ cells per regime---the taper regime is benign (median $1.8^\circ$,
ninetieth percentile $5.4^\circ$) because flux runs straight down a smoothly
tapering axis, whereas the bottleneck regime has a far heavier tail (median
$3.0^\circ$, ninetieth percentile $21.2^\circ$, a $3.95\times$ larger tail),
concentrated exactly at the neck where flux magnitude is highest and direction
matters most. The regime in which an eikonal ATR would be \emph{used to detect}
$\rhou>1$ is thus the regime in which it is \emph{least reliable}. Separately, a
sharp constriction produces a small flux-\emph{reversal} pocket just past the neck
(about $3.9\%$ of $\lvert J\rvert$ at the tightest necks), where the true current
briefly points back toward the interior and eikonal descent---which always points
toward $\Sig$---is about $178^\circ$ wrong. This is a localized re-entrant-corner
effect of classical elliptic type: it scales with corner sharpness and vanishes to
zero in a straight channel, so it is a categorical rather than an angular failure,
reported separately from the corridor error and not its driver. The consequences
for the protocol are concrete: an eikonal ATR is admitted to the contest only with
its corridor error and corner-reversal share reported per geometry, and it may not
be adopted for bottleneck prediction on the strength of taper accuracy. The
flux-tube diagnostic built from the true $J$, which has no such error by
construction, is the reference against which every cheap surrogate is scored.

\section{Two geometric regimes and the ratio that separates them}\label{sec:regimes}

Replacing the single ``narrowing raises resistance'' story with the corrected
mathematics produces a genuine dichotomy, and both halves are numerically real.
In a \emph{volume-graded taper}, downstream demand falls with cross-section fast
enough that the burden drops below the distance prediction; the power-law horns
realize it, with the tube quantity shrinking $0.167\to0.046$ as computed in
\eqref{eq:graded}. In a \emph{demand-loaded bottleneck}, a narrow section still
carries large downstream demand, and the burden can exceed a matched-distance
reference. Table~\ref{tab:regimes} records the two.

\begin{table}[t]
\centering
\begin{tabular}{@{}p{0.20\textwidth}p{0.34\textwidth}p{0.34\textwidth}@{}}
\toprule
\textbf{Regime} & \textbf{Mechanism} & \textbf{Effect and evidence}\\
\midrule
Volume-graded taper &
downstream demand falls with cross-section fast enough &
burden falls \emph{below} the distance prediction, $\rhou\ll1$; power-law horns,
tube quantity $0.167\to0.046$\\[4pt]
Demand-loaded bottleneck &
a narrow section still carries large downstream demand &
burden \emph{exceeds} a matched-distance reference, $\rhou>1$; a matched-Krogh
bottleneck attains $\rhou=1.83$\\
\bottomrule
\end{tabular}
\caption{The two regimes of the transport burden. The organizing quantity is
neither distance nor narrowness but the local demand-to-cross-section ratio
\eqref{eq:localratio}.}
\label{tab:regimes}
\end{table}

The quantity that organizes both is the \emph{local ratio}
\begin{equation}\label{eq:localratio}
\frac{\Qdn(s)}{P(s)\,\Acr(s)},
\end{equation}
the maintenance demand routed through each available transport cross-section. A
narrow section is cheap when little demand remains beyond it; a modest bottleneck
is severe when it must serve a large downstream region. This is a stronger and
more falsifiable statement than the original distance conjecture, and it is what
Phase~A is built to test. The existence of the $\rhou>1$ regime is not
hypothetical: a bottleneck matched to its Krogh reference on maximum distance,
total demand, and volume attains a true-field amplification $\rhou=1.83$. The
\emph{magnitude} of achievable amplification, and whether a closed-form
bottleneck family exists, remain the sharpest open targets of the program.

\section{Two amplification ratios, kept distinct}\label{sec:ratios}

An earlier draft wrote a single symbol $\rho$ for two different things, and the
resulting confusion is worth preventing permanently. The \emph{true-field
amplification} compares the exact burden to that of the matched Krogh reference
geometry,
\begin{equation}\label{eq:rhou}
\rhou=\frac{\norminf{u}}{u_{\mathrm{Krogh}}},
\end{equation}
so $\rhou>1$ means the geometry is more burdensome than its Krogh reference,
$\rhou<1$ less, $\rhou=1$ neither. Horns have $\rhou\ll1$; the demand-loaded
bottleneck has $\rhou>1$. The \emph{surrogate amplification} is the same
comparison computed from the deployable surrogate,
\begin{equation}\label{eq:rhohat}
\rhohat=\frac{R^{\text{surr}}_{\mathrm{ATR}}}{R^{\text{surr}}_{\mathrm{ATR,Krogh}}}.
\end{equation}
Phase~A asks precisely whether $\rhohat$ predicts $\rhou$. The two must be kept
apart exactly because the horn family lives at $\rhou\ll1$ while the bottleneck
family lives at $\rhou>1$: a single word ``amplification'' spanning both would be
actively misleading.

\section{The Phase-A surrogate contest}\label{sec:phaseA}

The scientific object of this paper is not a biological result but a protocol
whose first phase settles a mathematical question with no biology in it at all:
\emph{which} geometry-only summary, if any, tracks the exact torsion burden.

\paragraph{The four phases.} Phase~A runs the contest on synthetic domains matched
on maximum Euclidean distance, volume, total demand, diffusivity, and
supply-boundary geometry (Section~\ref{sec:supply}), with cross-sectional
restriction varied systematically across \emph{both} regimes. For each geometry we
compute and rank against the exact $\norminf{u}$ a ladder of candidates:
Euclidean maximum distance; geodesic distance; a skeleton or medial-axis ATR; an
eikonal ATR; the flux-tube diagnostic from the true $J$; and the exact field
itself. Phase~B applies the survivors to real segmented geometries and confirms
calibration against the full PDE on the same geometry---the gatekeeper again---
with no outcome variable yet. Phase~C introduces a matched biological readout and
tests incremental predictive power beyond distance and the standard covariates,
with the outcome-to-functional mapping preregistered
(Section~\ref{sec:outcome}). Phase~D asks the developmental question---do
organisms alter architecture in the predicted direction---and is reached only
after the physical predictor has passed; cross-sectional tissue data can show
association only, so the threshold-activation claim ultimately wants
developmental or perturbational evidence.

\paragraph{An ablation ladder, so no surrogate is a black box.} Reporting only a
surrogate's final score against $u$ would leave unknown \emph{which} ingredient
produced it. Every surrogate is therefore reported along a fixed ladder that adds
one information source at a time. The first rung uses Euclidean distance only; the
second replaces it with geodesic distance; the third adds the local cross-section,
i.e.\ restriction; the fourth instead adds downstream demand; the fifth is the
full surrogate, in which restriction and demand are combined multiplicatively and
integrated along the path; and the sixth rung is the exact PDE burden $u$, the
ceiling. For each rung we report the same package used everywhere in the
protocol---grouped cross-validated $R^2$, median relative error, Spearman
correlation, calibration slope, and matched-pair ordering---so that the
question ``which rung produces the gain'' has a quantitative answer that is itself
a result, independent of whether the final surrogate is ever adopted.

\paragraph{The ladder is decidable, and the decomposition is stark.} A ladder is
only meaningful if its rungs actually separate; two collinear adjacent rungs would
make the decomposition illusory. On a $16$-geometry family spanning both regimes
($8$ tapers and $8$ demand-loaded bottlenecks, all matched on maximum distance),
the largest adjacent-rung collinearity is $\lvert\text{corr}\rvert=0.369$---far
below the $\approx0.98$ that would signal redundancy, so the rungs carry
independent information. On that matched family the decomposition is decisive.
Because maximum distance is held fixed, Euclidean distance is blind
($\text{corr}=0.00$ with $u$); cross-sectional restriction alone is partial and
\emph{inversely} aligned ($-0.85$); downstream demand alone is only weakly aligned
($+0.23$); and only their path-integrated product tracks the true burden
($+0.92$). The predictive gain is neither geometry nor demand in isolation but the
restriction$\,\times\,$demand interaction---a concrete, pre-biological instance of
the thesis that burden tracks $\Qdn(s)/(P(s)\Acr(s))$. Whether this decomposition
survives on real segmented geometries is itself a Phase-B question.

Where a rung's path or direction is built from an eikonal field, its error is not
uniform across the family, for the reasons of Section~\ref{sec:tiers}: small in
the taper regime but with a roughly fourfold heavier tail in the bottleneck, plus
the re-entrant-corner reversal that eikonal cannot represent. Each
eikonal-dependent rung is therefore reported separately by regime, so that a rung
which looks adequate on the taper-dominated average is not credited in the
bottleneck regime it exists to detect.

\paragraph{Four models and the discipline of comparing them.} The biological test
of Phase~C is organized around four nested-or-not models,
\begin{align*}
M_0&:\ \text{outcome}\sim\text{distance},\\
M_1&:\ \text{outcome}\sim\text{distance}+\text{demand}+\text{diffusivity}
        +\text{volume}+\text{covariates},\\
M_2&:\ M_1+\RATR^{\text{surr}},\\
M_3&:\ M_1+\RG\quad(\text{the exact-}u\text{ summary}).
\end{align*}
The comparisons must respect nesting. $M_1$ versus $M_2$ and $M_1$ versus $M_3$
are nested and are compared by likelihood-ratio tests together with grouped-CV
$\Delta R^2$ and AIC/BIC. But $M_2$ versus $M_3$ compares a surrogate against the
exact burden and is \emph{generally non-nested}: it must be compared by grouped
cross-validation, information criteria, or Vuong-type tests, never by a likelihood
ratio. Crucially, ``$M_2\approx M_3$'' is not a discretionary judgment; we declare
\emph{surrogate parity} only when $M_2$'s grouped out-of-sample performance falls
within a preregistered tolerance of $M_3$---for instance $\Delta R^2_{\mathrm{CV}}
\le0.02$ with comparable calibration slope and matched-pair ordering. The
interpretation then separates two failure modes that a single test would confound:
if $M_3>M_1$ but $M_2\not>M_1$, the physical burden matters and the surrogate
failed; if neither beats $M_1$, modeled burden adds nothing in that system; and if
$M_2\approx M_3>M_1$, burden matters and the surrogate captures it---the target.

\paragraph{Why a high correlation is not enough.} A surrogate can correlate almost
perfectly with the truth and still invert the comparisons that matter. In a direct
construction, a surrogate with Pearson correlation $0.995$ against $u$ nonetheless
reverses the ordering of $20$ matched-distance pairs. The minimum acceptance
package is therefore not a correlation but the full set: cross-validated $R^2$,
median relative error, Spearman correlation, calibration slope and intercept, and
above all \emph{pairwise-ordering preservation at matched distance}, together with
out-of-sample performance and behavior under severe bottlenecks. The surrogate
must beat maximum distance on geometries constructed so that the two disagree---not
on a sample where distance already happens to work.

\section{Matching the summary to the outcome}\label{sec:outcome}

The scalar $\norminf{u}$ is the right summary for a specimen-level failure
threshold and the wrong one for a local or averaged outcome; the mapping from PDE
summary to biological readout must be fixed before Phase~C, not chosen to fit.
Table~\ref{tab:outcome} records the preregistered correspondence.

\begin{table}[t]
\centering
\begin{tabular}{@{}p{0.42\textwidth}p{0.46\textwidth}@{}}
\toprule
\textbf{Biological outcome} & \textbf{Matched functional of $u$}\\
\midrule
local osteocyte viability & $u(x)$, or a local-neighborhood summary\\
total necrotic fraction & the measure of $\{x:u(x)>\Delta C_{\text{avail}}\}$\\
mean oxygenation & a spatial or demand-weighted mean of $u$\\
architectural density & a percentile, or integrated exposure
$\int\mathbf{1}[u>\theta]$\\
\bottomrule
\end{tabular}
\caption{Preregistered mapping from the PDE summary to the biological outcome.
Choosing the summary after seeing the data would invalidate the incremental-
predictivity test of Phase~C.}
\label{tab:outcome}
\end{table}

\section{Supply-boundary geometry as a matched variable}\label{sec:supply}

A clean bottleneck experiment must hold fixed not only the maximum distance but
the effective supply interface, so that interior cross-sectional restriction is
the sole varied factor. Phase~A therefore controls or explicitly varies the
supply set $\Sig$: its total Dirichlet surface area, the number of disconnected
supply patches, their placement, the boundary concentration, and---if a Robin
contact condition later replaces ideal Dirichlet---the contact conductance. The
matched-Krogh bottleneck that attains $\rhou=1.83$ was matched to its reference on
maximum distance, total demand, and volume for exactly this reason; without the
supply interface held fixed, an apparent amplification could be an artifact of a
changed boundary rather than of interior geometry.

\section{A dimensionless maintenance number}\label{sec:pimaint}

With the unit split of Section~\ref{sec:forced} in hand, the feasibility of
maintenance admits a clean dimensionless form. Writing $C_\Sig$ for the supply
concentration and $C_{\text{crit}}$ for the viability threshold, set
\begin{equation}\label{eq:pimaint}
\Pimaint=\frac{\dCtube}{C_\Sig-C_{\text{crit}}},\qquad\text{and, for uniform }q,\ \
\dCtube=q\,\Rgeom .
\end{equation}
Then $\Pimaint<1$ is viable, $\Pimaint\approx1$ is threshold, and $\Pimaint>1$
predicts depletion beyond viability. The correction that makes this honest is
dividing by the \emph{margin above threshold} $C_\Sig-C_{\text{crit}}$, not by
$C_\Sig$ alone. One should report $\Pimaint$ alongside the true-field
amplification $\rhou$ of \eqref{eq:rhou}, and never alongside the comparison-field
quantity, which shrinks with taper and would understate the load.

\section{Branch points: competing primary models, not a footnote}\label{sec:branch}

At a branch point the upstream demand must be split among daughter paths in order
to evaluate $\Qdn(s)/\Acr(s)$, and there are two natural rules. An
\emph{area-proportional} split sends demand in proportion to the downstream
cross-section carried by each daughter, in the spirit of Murray's law and a
Kirchhoff balance~\cite{murray1926}; a \emph{demand-proportional} split sends it
in proportion to the downstream integrated demand of each subtree. In the real
PDE the branch flux is set jointly by conductance, geometry, sources, and
gradients, and neither shortcut is automatically the biological rule. We therefore
run both as \emph{competing primary models}, not as a sensitivity appendix, and
report their full model comparisons side by side. If the ATR surrogate wins under
both rules, the result is robust and partition-independent; if it wins under only
one, that is informative about \emph{what} the functional detects---whether
demand-weighting or area-weighting carries the signal---rather than a nuisance to
be hidden.

\section{The lacuno-canalicular test case and its controls}\label{sec:lcn}

The lacuno-canalicular network (LCN) of bone is the strongest biological
discriminator available, because its measured, restricted cross-section departs
sharply from Krogh distance; it is exactly where distance and ATR most disagree.
For that same reason it must \emph{not} be the first implementation. An LCN result
carries meaning only after the surrogate has passed the gatekeeper on controlled
synthetic horns and branching networks where the true $u$ is known.

Even then, a specimen-level LCN test discharges only under conditions that the
protocol states in advance. The supply set may exceed the Haversian canal, drawing
on Volkmann canals, neighboring osteons, bone and marrow surfaces, and pressure
gradients. Transport may include load-induced advection, requiring
$-\nabla\!\cdot(P\nabla c)+\mathbf{v}\!\cdot\!\nabla c=q$ rather than pure
torsion. Canalicular density cannot be both predictor and outcome---it constructs
the cross-section $\Acr$, so the dependent variable must be independent of it
(osteocyte viability, lacunar occupancy, recovery time, solute penetration, or
metabolic markers). And radial averages erase the very bottlenecks, connectivity,
and dead-ends that the functional is meant to detect, so a confirmatory test needs
specimen-level three-dimensional geometry paired with matched specimen-level
readouts, not a mean profile paired with unrelated viability data. Passing all
four, cross-sectional data still show association only; the threshold-activation
claim wants the developmental or perturbational evidence of Phase~D.

Two controls frame the discriminator. \emph{Spheroids} are a principled
near-Krogh negative control, with $\rhou\approx1$ expected; a null result there
confirms the domain of relevance rather than refuting the concept.
\emph{Muscle} is a stress test rather than a confirmation: myoglobin-facilitated
transport and the mitochondrial reticulum can negate classical diffusion
distances, so both ATR and Krogh distance may fail there, and this should be
flagged prominently because muscle-literate readers will probe it.

\section{Interpretation: an accessibility potential, and the demotion of the
free-energy reading}\label{sec:interp}

The empirical content of this program stands with the free-energy principle (FEP)
discarded. FEP is, at most, an optional and terminal interface language for
maintenance under uncertainty; it does not generate the transport functional, and
the explanatory arrow stays physical: geometry sets the PDE burden $\RG=u/\Qstar$,
which sets maintenance feasibility. We state this demotion plainly so that no
predictive weight rests on FEP.

What survives, and is measurable, is a reading of $u$ as an \emph{accessibility-
cost potential}: the scalar cost of maintaining an interior point under a supply
boundary, distributed demand, and constrained transport geometry, with $\RATR$ its
pathwise accumulation. The relation is clean---$\RATR$ rising is accessibility
falling. When passive accessibility becomes too costly, a system can change
geometry, widen cross-section, alter permeability, add sources, or grow active-
maintenance architecture. This is a material realization of a single principle,
offered as interpretation of the same physics rather than as an independent source
of predictions:
\begin{quote}
Form does not respond to raw distance. It responds to the accumulated difficulty
of maintaining relation across its own geometry.
\end{quote}
Unlike the FEP bridge, this reading asserts no upstream awareness or optimization;
it supplies a middle layer from relational geometry to biological architecture
that is, in principle, measurable.

\section{Status of the claims}\label{sec:status}

Because the protocol mixes theorems, modeling choices, and open questions, we
record which is which without decoration.

\emph{Forced by the current mathematics.} The Green representation
\eqref{eq:green} holds where the mixed Green function exists (a.e.\ $x$, pointwise
under Phase-A regularity). Distance to the supply set yields no universal lower
bound. The explicit horn comparison field \eqref{eq:vfield} solves the PDE with
exactly vanishing lateral-wall flux, and its apex obeys the graded identity
\eqref{eq:graded}; the field gives a rigorous upper bound \eqref{eq:upper} on the
true burden. A distance-only burden can overestimate the true torsion by an
arbitrarily large factor on the horn family (failure factor $6\to22$). For uniform
demand, $\dCtube=q\,\Rgeom$ \eqref{eq:qrel}. A demand-loaded bottleneck can be more
burdensome than its matched Krogh reference, $\rhou=1.83$, establishing the
$\rhou>1$ regime. On a matched-distance family spanning both regimes the ablation
rungs separate and the predictive gain is the restriction$\,\times\,$demand
interaction, not distance, restriction, or demand alone (correlations with $u$ of
$0.00$, $-0.85$, $+0.23$, $+0.92$). The eikonal-surrogate error is regime-
dependent---small in the taper regime, with a roughly fourfold heavier tail in the
bottleneck, plus a re-entrant-corner flux-reversal pocket the eikonal cannot
represent---so it is least reliable in the very regime it would be used to detect.

\emph{Declared constructions.} The distinct objects $\Rgeom$ and $\dCtube$; the
normalized burden $\RG=u/\Qstar$; the flux-tube diagnostics; the eikonal,
skeleton, and geometry-only surrogates; the area- and demand-proportional branch
partitions; the maintenance number $\Pimaint$; the amplifications $\rhou$ and
$\rhohat$; and the accessibility-potential reading are all modeling choices, sound
but not forced.

\emph{Open.} Whether the true horn apex $u(0)$ equals the graded tube integral
rather than merely being bounded by it; the magnitude of achievable $\rhou>1$ and
any closed-form bottleneck family; whether a uniquely preferred general-domain
surrogate exists; whether ATR beats distance plus conventional covariates on real
biology; whether a universal persistence envelope exists; and any FEP derivation.

\section{Computational verification}\label{sec:verify}

Every displayed algebraic identity and rational inequality above, and every
numerical claim about the two regimes, is machine-checked by the companion
harness \texttt{surrogate\_contest.py} (exit~$0$), which builds each synthetic
geometry, solves \eqref{eq:pde} to obtain the exact $u$, and scores the
candidates. The eight check families are: (C1) eikonal descent differs from the
true flux, $10.6^\circ$, with eikonal $0^\circ$ from the supply direction; (C2) a
matched-distance, matched-demand bottleneck raises the true $\norminf{u}$ by
$1.66\times$, tracked by the exact tube reduction at $1.63\times$---distinct
ratios---and missed by geodesic distance ($1.6\%$); (C3) a Pearson-$0.995$
surrogate inverts $20$ matched pairs; (C4) the horn comparison quantity shrinks
with taper ($0.167\to0.046$) while the distance-only failure factor grows
($6\to22$); (C5) the uniform-demand reduction $\dCtube=q\,\Rgeom$ holds with no
double-count; (C6) a demand-loaded bottleneck attains $\rhou=1.83>1$ against its
matched Krogh reference---the open regime, numerically realized; (C7) the ablation
ladder is decidable (largest adjacent-rung $\lvert\text{corr}\rvert=0.369$) and
mechanism-isolating (the gain is the restriction$\,\times\,$demand interaction,
with correlations $0.00/-0.85/+0.23/+0.92$); and (C8) the eikonal-surrogate error
is regime-dependent (taper corridor ninetieth percentile $5.4^\circ$ versus
bottleneck $21.2^\circ$) with a small re-entrant-corner reversal pocket at sharp
necks. The corner-reversal effect was probed for standalone significance by a
second companion script, \texttt{reversal\_generality.py}: it is generic and
resolution-independent but a local elliptic-corner phenomenon that scales with
neck sharpness and vanishes in a straight channel---classical, not novel---and is
folded in as a caveat rather than a result.

\section{Outlook}\label{sec:outlook}

The logical spine of the program is short and depends on no equality the horn
paper did not establish:
\begin{center}
exact torsion field $u=\int G\,q$\ \ $\to$\ \ comparison-field horn identity
$v(0,0)=\tfrac{q}{P}\int V/\Acr$ (an upper bound on $u$)\ \ $\to$\ \ geometry-only
surrogate contest across both regimes (Phase~A)\ \ $\to$\ \ real-geometry PDE
calibration, the gatekeeper (Phase~B)\ \ $\to$\ \ biological-outcome test with the
four models and both branch partitions (Phase~C)\ \ $\to$\ \ a possible
persistence principle, developmentally probed (Phase~D).
\end{center}
The highest-value next step is Phase~A: to determine which general-domain
surrogate, if any, tracks the exact torsion burden across both the taper and the
demand-loaded-bottleneck regimes under adversarial perturbation---with the
bottleneck family, now numerically real at $\rhou=1.83$, as the sharpest probe.

\begin{thebibliography}{9}
\bibitem{torsion} Echo-S Research,
\emph{Forced Mechanics of the Mixed-Boundary Torsion Problem: Exact Laws, Sharp
Star-Shaped Bounds, and the Closure of the Distance Conjecture}, Echo S Research
archive, 2026. (Companion paper; ``the horn paper'' of this text.)
\bibitem{krogh1919} A.~Krogh,
\emph{The number and distribution of capillaries in muscles with calculations of
the oxygen pressure head necessary for supplying the tissue}, J.~Physiol.\
\textbf{52} (1919), 409--415.
\bibitem{murray1926} C.~D.~Murray,
\emph{The physiological principle of minimum work applied to the angle of
branching of arteries}, J.~Gen.~Physiol.\ \textbf{9} (1926), 835--841.
\bibitem{friston2010} K.~Friston,
\emph{The free-energy principle: a unified brain theory?}, Nature Reviews
Neuroscience \textbf{11} (2010), 127--138. (Cited only to be demoted; see
Section~\ref{sec:interp}.)
\end{thebibliography}

\end{document}
